% Plan view of the Case-C direction rule, seen along the inward surface normal.
% One panel per commanded rotation direction. Each panel carries the surface
% tangents, the signed rotation axis u_rot, and the tangential lever selected by
% eq:experimental_general_lever_rule. The lever is perpendicular to the
% rotation axis, as recorded by the right-angle marker at the origin.
%
% The four rotation axes lie at 115, 70, 25 and 160 degrees from t_1, and the
% selected levers therefore lie at 25, -20, -65 and 70 degrees. Those are the
% directions of the lever vectors listed in tab:experimental_case_c_vectors.
%
% Uses only tikz + arrows.meta, which config/packages.tex loads. Colour follows
% the established meaning: orange for what a parameter selects (the lever),
% grey for the frame and the commanded rotation axis.
\begin{tikzpicture}[
    font=\small,
    ax/.style={-{Latex[length=1.8mm]}, gray!60!black},
    rotax/.style={{Latex[length=1.8mm]}-{Latex[length=1.8mm]}, densely dashed,
                   gray!55!black},
    lev/.style={-{Latex[length=2.2mm]}, thick, orange!75!black},
    lbl/.style={font=\footnotesize, inner sep=1.8pt},
  ]

\foreach \dx/\ua/\name in {%
    0/115/{about $y_{\mathrm{EE}}$},
    3.4/70/{$-45^\circ$ diagonal},
    6.8/25/{about $x_{\mathrm{EE}}$},
    10.2/160/{$+45^\circ$ diagonal}}
{
  \begin{scope}[xshift=\dx cm]
    \pgfmathsetmacro{\ra}{\ua-90}

    % ---- surface tangents ----------------------------------------------
    \draw[ax] (0,0) -- (0.95,0) node[lbl, right] {$t_1$};
    \draw[ax] (0,0) -- (0,0.95) node[lbl, above] {$t_2$};

    % ---- commanded rotation axis, drawn through the origin in both senses
    % The label sits on the far end so it does not crowd the lever label of
    % the neighbouring panel.
    \draw[rotax] (\ua+180:1.08) -- (\ua:1.08);
    \node[lbl, gray!55!black] at (\ua+180:1.32) {$u_{\mathrm{rot}}$};

    % ---- selected tangential lever --------------------------------------
    \draw[lev] (0,0) -- (\ra:1.08);
    \node[lbl, orange!65!black] at (\ra:1.30) {$r_{c,t}$};

    % ---- right angle between the two ------------------------------------
    \begin{scope}[rotate=\ra]
      \draw[gray!60, thin] (0.17,0) -- (0.17,0.17) -- (0,0.17);
    \end{scope}

    \node[lbl] at (0,-1.62) {\name};
  \end{scope}
}

\end{tikzpicture}
